\section{Chapitre~\ref{sec:dishbrain}. DishBrain}

La conception du banc et les résultats du jeu proviennent d'une seule étude expérimentale et de sa suite; les formules du bruit et de la dérive vers les bons réglages sont établies dans le chapitre et vérifiées sur le modèle du livre, tandis que le mécanisme d'apprentissage dans la boîte n'est pas établi.

{\footnotesize
\setlength{\tabcolsep}{3pt}\renewcommand{\arraystretch}{1.15}
\begin{longtable}{>{\raggedright}p{0.5cm}>{\raggedright}p{4.0cm}>{\raggedright}p{5.6cm}>{\raggedright}p{2.3cm}>{\raggedright\arraybackslash}p{1.7cm}}
\toprule
N\textsuperscript{o} & Affirmation & Expérience et ce qui est mesuré & Source & Statut \\
\midrule\endfirsthead
\toprule
N\textsuperscript{o} & Affirmation & Expérience et ce qui est mesuré & Source & Statut \\
\midrule\endhead
1 & Banc: environ $800\,000$ cellules de cortex de souris ou environ $10^6$ cellules humaines par puce; $26\,000$ électrodes sur $8$~mm$^2$, enregistrement jusqu'à $1024$, stimulation par $8$ & Méthodes de l'étude: puce MaxOne, $26\,000$ électrodes de platine sur $8$~mm$^2$, enregistrement de $1024$ canaux à $20\,000$ échantillons par seconde; techniquement, on peut stimuler jusqu'à $32$ électrodes, $8$ ont pu l'être indépendamment; les cultures de souris ont été utilisées à partir d'environ $10$~jours & \cite{Kagan2022} & mesuré \\
2 & Boucle cadencée à $10$~ms: les impulsions des zones motrices déplacent la raquette (fig.~\ref{fig:dishloop}) & Les impulsions sont comptées par fenêtres de $10$~ms; le délai entre une impulsion et le stimulus de réponse, d'environ $5$~ms, est fixé par le tampon du matériel & \cite{Kagan2022} & mesuré \\
3 & Entrée: code de place (laquelle des $8$ électrodes) et fréquence de $4$ à $40$~imp./s & Dans l'expérience principale, la fréquence de stimulation croît linéairement de $4$~Hz quand la balle est près du mur du fond à $40$~Hz près de la raquette; l'amplitude des impulsions de la balle est de $75$~mV; les impulsions sont rectangulaires biphasiques & \cite{Kagan2022} & mesuré \\
4 & Sortie $\Delta p=c\,(n_\uparrow-n_\downarrow)$~\eqref{eq:dishreadout}, bruit du compte par fenêtre $1/\sqrt{RT}$ & La règle de différence entre deux zones est décrite dans l'étude (avec des poids des zones selon leur activité), la formule exacte n'est pas donnée. L'estimation du bruit découle du comptage poissonien et est établie dans le chapitre & \cite{Kagan2022}; démonstration dans le chapitre & théorie \\
5 & Professeur: après un renvoi, stimulus prévisible de $75$~mV, $100$~imp./s, $0.1$~s; après un raté, stimulus imprévisible de $150$~mV, $5$~imp./s, $4$~s en des lieux et instants aléatoires, puis $4$~s de pause & Méthodes: après un renvoi, $75$~mV, $100$~Hz, $100$~ms sur les $8$ électrodes à la fois; après un raté, $150$~mV, $5$~Hz sur des électrodes aléatoires à des instants aléatoires pendant $4$~s, puis $4$~s sans stimulation. La culture ne contient pas de dopamine & \cite{Kagan2022} & mesuré \\
6 & Le réseau agit de façon à rendre son entrée prévisible & Le principe de l'énergie libre est une proposition théorique dont les auteurs ont tiré le protocole; l'expérience est compatible avec lui, mais ne le distingue pas de mécanismes plus simples (lignes~7--8) & \cite{Friston2010,Kagan2022} & hypothèse \\
7 & Si l'échec secoue le réseau et non le succès, le temps passé dans un réglage vaut $\rho\propto1/P_{\text{raté}}$~\eqref{eq:dishdensity} (proposition~\ref{prop:dishscramble}) & Découle de l'équilibre des flux dans une chaîne de Markov à secousses symétriques, établi dans le chapitre. Aucun test direct sur une culture & démonstration dans le chapitre & théorie \\
8 & Le modèle du \og brassage après un raté\fg{} apprend: proportion de renvois $0.21\to0.38$; avec une secousse après chaque échange $\approx0.12$, sans secousse $0.19$ (fig.~\ref{fig:dishmodel}) & Calcul, script \texttt{models/dishbrain\_model.py}, $40$ cultures modèles de $2000$ échanges chacune: $0.21\to0.38$; $0.13\to0.11$; $0.19\to0.19$ & --- & modèle \\
9 & Les séries de renvois ne s'allongent que chez les cultures vivantes en boucle complète; les contrôles n'apprennent pas & $399$ sessions de $20$~min: $9$ cultures de souris ($101$ sessions), $11$ humaines ($138$), $6$ puces avec milieu seul ($80$), $20$ cultures au repos ($42$), $3$ simulations ($38$). La longueur moyenne d'un échange sur les $15$ dernières minutes n'est supérieure à celle des $5$ premières que chez les cultures vivantes; les cellules non neuronales (HEK293T) et les puces avec milieu seul ne montrent aucun effet & \cite{Kagan2022} & mesuré \\
10 & Hausse significative en une session seulement avec rétroaction complète; avec du silence à la place du stimulus, le jeu est meilleur en fin de session que sans rétroaction, mais l'effet est moindre & $486$ sessions sur $3$ jours: \og stimulus\fg{} ($n=27$), \og silence\fg{} ($17$), \og sans rétroaction\fg{} ($15$). La hausse au fil du temps n'est significative que dans la condition \og stimulus\fg{}; en fin de session, la condition \og silence\fg{} dépassait \og sans rétroaction\fg{} avec un effet moindre & \cite{Kagan2022} & mesuré \\
11 & L'effet est faible; savoir si le réseau apprend durablement reste une question ouverte & Au sein d'une session, l'amélioration est robuste, mais, de l'aveu même des auteurs, aucun apprentissage stable d'une session à l'autre n'a été observé sur plusieurs jours & \cite{Kagan2022} & mesuré \\
12 & Pendant le jeu, le réseau se rapproche du régime critique, et plus il s'en approche, meilleur est le jeu & Analyse des avalanches d'activité des mêmes cultures: une entrée structurée rapproche le réseau du régime critique, et la qualité du jeu est corrélée à cette proximité; mais la criticité seule, sans information sur les conséquences des actions, ne suffit pas à l'apprentissage & \cite{Habibollahi2023} & mesuré \\
13 & La \og sensibilité\fg{} de la culture n'est pas démontrée & Article de réponse de trente chercheurs: un changement de comportement en boucle fermée ne prouve ni intelligence ni sensations; aucune expérience ne le démontre & \cite{Balci2023} & hypothèse \\
14 & Quatrième contrôle proposé: un stimulus prévisible après un raté, de même durée et de même énergie (section~\ref{sec:dishspec}) & Cette expérience n'a pas été réalisée; c'est une proposition du livre & --- & hypothèse \\
\bottomrule
\end{longtable}}
